\clearpage
\section{ファイルリスト（写真の読み込みと管理）}\label{sec:filelist}

\begin{figure}[H]
\centering
\begin{screenshot}[0.36\linewidth]{filelist}
  \calloutright{1}{reset}
  \calloutleft{2}{base}
  \calloutleft{3}{types}
  \calloutleft{4}{add}
  \calloutright{5}{clear}
  \calloutleft{6}{star}
  \calloutright{7}{remove}
  \calloutleft{8}{unstar}
\end{screenshot}
\caption{ファイルリスト}
\end{figure}

\begin{callouts}
\co{1} \ui{すべてクリア}：読み込んだ写真・合成の結果・マスク・設定をすべて消して、起動した直後の状態に戻します。確認のメッセージが出ます。取り消せないので注意してください。
\co{2} \textbf{基準画像}：いまの基準画像の名前と撮影情報です。
\co{3} \textbf{種類の切り替え}：\ui{Light}・\ui{Dark}・\ui{Flat}・\ui{Bias}を切り替えます。かっこの中は読み込んだ枚数です。\ui{＋ 追加}やドラッグ＆ドロップで読み込むと、いま選んでいる種類に入ります。
\co{4} \ui{＋ 追加}：写真を選んで、いま選んでいる種類に読み込みます。
\co{5} \ui{リストをクリア}：いま選んでいる種類（Light なら Light だけ）の写真をすべて外します。
\co{6} 黄色の\ui{★}が基準画像です。いま表示している写真は、左に青い帯が付き、名前が太字になります。
\co{7} \ui{×}：その写真をリストから外します（ファイルそのものは消えません）。
\co{8} \ui{☆}：押すと、その写真が基準画像になります。
\end{callouts}

写真の名前の下には、撮影に使ったレンズや、焦点距離・F値・ISO・露出時間が表示されます。
写真の行をクリックすると、その写真がプレビューに表示されます。

\begin{memo}[取り消し]
写真を追加したり外したりした操作は、\menu{編集 → 取り消す}（\key{⌘}\,\key{Z}）で1つずつ戻せます。
\end{memo}

\subsection{Dark・Flat・Bias（補正用の写真）}\label{sec:calibration}

Dark・Flat・Bias は、センサーやレンズのくせを取り除くための写真です（\pageref{sec:requirements}ページの表も見てください）。
なくても合成できます。使う場合は、種類を切り替えてから\ui{＋ 追加}で読み込むだけで、合成のときに自動で補正します。

\begin{center}
\small
\begin{tabularx}{\linewidth}{l X X}
\toprule
種類 & 撮り方 & 取り除けるもの \\
\midrule
Dark & レンズキャップをして、Light と同じ露出時間・ISO で、できれば同じ気温のときに数枚〜十数枚 & 熱によるノイズ、ホットピクセル（明るい点）、アンプグロー（画面の端の明るいかぶり） \\
Flat & Light と同じレンズ・F値・ピントで、明るさが均一な面を数枚 & 周辺減光（四隅が暗くなる）、センサーのごみの影 \\
Bias & レンズキャップをして、いちばん短い露出時間・Light と同じ ISO で数枚〜十数枚 & センサーの基準の明るさのむら \\
\bottomrule
\end{tabularx}
\end{center}

\begin{hint}
何枚か読み込むと、自動で1枚にまとめて（中央値）から使います。Dark を読み込んだ場合は、Bias はなくてもかまいません。
\end{hint}

\subsection{地上固定フレームの欄}

新星景モードで「地上固定フレーム」を登録すると、Light タブのいちばん上に\ui{地上固定フレーム}の欄ができ、登録した写真が並びます（\ref{sec:groundfixed}節）。
